\section*{Checklist}

\begin{enumerate}
\item For all models and algorithms presented, check if you include:
  \begin{enumerate}
  \item A clear description of the mathematical setting, assumptions, algorithm, and/or model. [Yes/No/Not Applicable]\\
  \textbf{Yes.} The main paper defines the offline MDP setting, TD3+BC backbone, BAR actor chain, routing, direct two-actor control, and the distinction between the ideal proximal map and the live one-Adam-step implementation.
  \item An analysis of the properties and complexity (time, space, sample size) of any algorithm. [Yes/No/Not Applicable]\\
  \textbf{Yes.} BAR stores $K$ persistent actors and makes $K$ actor optimizer calls per delayed update. The paper reports both the $O(K)$ actor-work scaling and one-stack H200 timing/memory measurements for $K=1$ and $K=4$.
  \item (Optional) Anonymized source code, with specification of all dependencies, including external libraries. [Yes/No/Not Applicable]\\
  \textbf{Partial.} The release contains BAR, the two-actor control, diagnostic runners/verifiers, compact scientific artifacts, and dependency files. Some historical runs predate the current provenance harness and therefore lack a complete executable historical snapshot.
  \end{enumerate}

\item For any theoretical claim, check if you include:
  \begin{enumerate}
  \item Statements of the full set of assumptions of all theoretical results. [Yes/No/Not Applicable]\\
  \textbf{Yes.} The first-hop target-perturbation result states its Lipschitz assumption; the Wasserstein/proximal and local explicit--implicit statements are explicitly restricted to deterministic policies, fixed critics/scales, and local conditions where appropriate.
  \item Complete proofs of all theoretical results. [Yes/No/Not Applicable]\\
  \textbf{Yes.} The supplementary material gives the behavior-anchor decomposition, inexact learned-critic margin, target-perturbation proof, and local explicit--implicit expansion.
  \item Clear explanations of any assumptions. [Yes/No/Not Applicable]\\
  \textbf{Yes.} The paper repeatedly distinguishes ideal fixed-critic statements from the persistent neural actor chain and reports a learned-ReLU activation-region residence audit to delimit the local interpretation.
  \end{enumerate}

\item For all figures and tables that present empirical results, check if you include:
  \begin{enumerate}
  \item The code, data, and instructions needed to reproduce the main experimental results (either in the supplemental material or as a URL). [Yes/No/Not Applicable]\\
  \textbf{Yes.} The repository contains fixed score matrices, compact diagnostic CSV/JSON artifacts, analysis/verifier code, experiment runbooks, and manuscript build instructions. Public D4RL datasets are cited.
  \item All the training details (e.g., data splits, hyperparameters, how they were chosen). [Yes/No/Not Applicable]\\
  \textbf{No.} All recoverable algorithmic settings are documented, but exact historical per-run manifests for proximal seeds 2--3 were not retained, and the exact historical Git revision of the earlier P3 two-actor control was never persisted. New locked experiments retain substantially stronger provenance.
  \item A clear definition of the specific measure or statistics and error bars (e.g., with respect to the random seed after running experiments multiple times). [Yes/No/Not Applicable]\\
  \textbf{Yes.} The paper defines the four-seed score grids, task-equal high-budget estimand, raw and seed-mean collapse counts, task-resampling intervals, paired P0 contrast, movement ratios, target-value RMS diagnostic, and the resampling/generalization scope.
  \item A description of the computing infrastructure used. [Yes/No/Not Applicable]\\
  \textbf{Partial.} The measured actor-cost audit records an H200/JAX 0.10.2 stack and new locked experiments record host/accelerator inventories. A complete independently verified hardware inventory for every historical headline run was not retained.
  \end{enumerate}

\item If you are using existing assets (e.g., code, data, models) or curating/releasing new assets, check if you include:
  \begin{enumerate}
  \item Citations of the creator if your work uses existing assets. [Yes/No/Not Applicable]\\
  \textbf{Yes.} D4RL, TD3+BC, CORL, and the relevant prior methods are cited.
  \item The license information of the assets, if applicable. [Yes/No/Not Applicable]\\
  \textbf{Yes.} The released repository identifies its license; upstream packages and benchmark assets retain their respective licenses.
  \item New assets either in the supplemental material or as a URL, if applicable. [Yes/No/Not Applicable]\\
  \textbf{Yes.} Compact P0--P3 evidence, verification receipts where available, and analysis code accompany the source release.
  \item Information about consent from data providers/curators. [Yes/No/Not Applicable]\\
  \textbf{Not Applicable.} The work uses public simulated-control benchmark data and no human-subject data.
  \item Discussion of sensitive content if applicable. [Yes/No/Not Applicable]\\
  \textbf{Not Applicable.} The assets contain simulated locomotion trajectories and no personal or sensitive information.
  \end{enumerate}

\item If you used crowdsourcing or conducted research with human subjects, check if you include:
  \begin{enumerate}
  \item The full text of instructions given to participants and screenshots. [Yes/No/Not Applicable]\\
  \textbf{Not Applicable.}
  \item Descriptions of potential participant risks, with links to Institutional Review Board approvals if applicable. [Yes/No/Not Applicable]\\
  \textbf{Not Applicable.}
  \item The estimated hourly wage paid to participants and the total amount spent on participant compensation. [Yes/No/Not Applicable]\\
  \textbf{Not Applicable.}
  \end{enumerate}
\end{enumerate}
